% ============ §32 The formal fibre (原书 pp.255-260) ============
\section{The formal fibre}\label{sec:32}\index{formal fibre, section 32}\index{regular homomorphism, section 32}\index{G-ring, section 32}\index{excellent ring, section 32}\index{quasi-excellent ring, section 32}

Let $(A, m)$ be a Noetherian local ring and $A^*$ its completion. The fibre ring of
the natural homomorphism $A \to A^*$ over any $p \in \Spec A$ is called a \emph{formal
fibre} of $A$ (although strictly speaking we should distinguish between the fibre
and the fibre ring, we will not do so in what follows). If $I$ is an ideal of $A$
then $(A/I)^* = A^*/IA^*$, so that a formal fibre of $A/I$ is also a formal fibre of
$A$.

Let $A$ be a Noetherian ring and $k \subset A$ a subfield. We say that $A$ is
\emph{geometrically regular} over $k$ if $A \otimes_k k'$ is a regular ring for every finite
extension $k'$ of $k$ (see \S\ref{sec:28}). This is equivalent to saying that $A_p$ is
geometrically regular over $k$ for every maximal ideal $p$ of $A$.

We say that a homomorphism $\varphi\colon A \to B$ of Noetherian rings is \emph{regular} if
$\varphi$ is flat, and for every $p \in \Spec A$, the fibre $B \otimes_A \kappa(p)$ of $\varphi$ over $p$
is geometrically regular over the field $\kappa(p)$.

A Noetherian ring $A$ is said to be a \emph{G-ring} (here G stands for Grothendieck)
if $A_p \to (A_p)^*$ is regular for every prime ideal $p$ of $A$; this means that all the
formal fibres of all the local rings of $A$ are geometrically regular.

\begin{theorem}\label{thm:32.1}
Let $A \xrightarrow{\varphi} B \xrightarrow{\psi} C$ be homomorphisms of Noetherian rings;
then
\begin{enumerate}[(i)]
  \item if $\varphi$ and $\psi$ are regular then so is $\psi\varphi$;
  \item if $\psi\varphi$ is regular and $\psi$ is faithfully flat then $\varphi$ is also regular.
\end{enumerate}
\end{theorem}

\begin{proof}
(i) Clearly $\psi\varphi$ is flat. For $p \in \Spec A$, write $K = \kappa(p)$, and let $L$ be a
finite extension field of $K$. Set $B_L = B \otimes_A L$, and $C_L = C \otimes_A L$; then the
homomorphism $\psi_L\colon B_L \to C_L$ induced by $\psi$ is also regular. Indeed, if $P$ is a
prime ideal of $B_L$ then $C_L \otimes_{B_L} \kappa(P) = C \otimes_B (B \otimes_A L) = C \otimes_A
\kappa(P) = C_{\kappa(P)}$, and hence if $F$ is a finite extension of $\kappa(P)$ then
$C_L \otimes_{B_L} F = C \otimes_A F$; setting $P \cap B = Q$, since $B_L$ is a finite $B$-module we
have $[\kappa(P) : \kappa(Q)] < \infty$, and hence $[F : \kappa(Q)] < \infty$, so that
$C \otimes_A F$ is a regular ring. Now $\varphi$ is regular, so that $B_L$ is a regular ring,
and hence by Theorem~\hyperref[thm:23.7]{23.7}, (ii), $C_L$ is a regular ring.

(ii) The flatness of $\varphi$ is obvious. If we let $p$, $K$ and $L$ be as above, then
$C_L$ is a regular ring and is flat over $B_L$, so that by Theorem~\hyperref[thm:23.7]{23.7}, (i), $B_L$
is also regular.
\end{proof}

\begin{theorem}\label{thm:32.2}
Let $\varphi\colon A \to B$ be a homomorphism of Noetherian rings, and assume that
$\varphi$ is faithfully flat and regular.
\begin{enumerate}[(i)]
  \item $A$ is regular (or normal, reduced, CM, or Gorenstein) if and only if $B$ has
        the same property.
  \item If $B$ is a G-ring then so is $A$ (the converse is not true).
\end{enumerate}
\end{theorem}

\begin{proof}
(i) follows from Theorem~\hyperref[thm:23.7]{23.7}, the corollaries to
Theorems~\hyperref[thm:23.9]{23.9} and \hyperref[thm:23.3]{23.3}, and Theorem~\hyperref[thm:23.4]{23.4}.

(ii) Let $p \in \Spec A$, choose $P \in \Spec B$ lying over $p$, and consider the
commutative diagram
\[
\begin{tikzcd}
(A_p)^* \arrow[r, "f^*"] & (B_P)^* \\
A_p \arrow[u, "\alpha"] \arrow[r, "f"] & B_P \arrow[u, "\beta"]
\end{tikzcd}
\]
Here $f$ is the map induced by $\varphi$, and $f^*$ is the map induced by $f$; and the
vertical arrows are the natural maps. Now $f$ and $\beta$ are both regular, and $f^*$ is
faithfully flat, so that according to the previous theorem, $\alpha$ is also regular.

To construct an example where $A$ is a G-ring and $B$ is not, we let $A = k$ be a
perfect field, and $B$ a regular local ring containing $k$. Then $k \to B$ is certainly
faithfully flat and regular, and $k$ is a field, and so trivially a G-ring. However,
there are known examples in which $B$ is not a G-ring. (See the appendix to [N1];
a counter-example is provided by the ring $R$ in (E3.1) if $\chark k = p$, and by $R$ in
Example 7 if $\chark k = 0$. In (E3.1) the field $k$ is not perfect, but $R$ is
geometrically regular over $k$.)
\end{proof}

\begin{theorem}\label{thm:32.3}
A complete Noetherian local ring is a G-ring.
\end{theorem}

\begin{proof}
Let $A$ be a complete Noetherian local ring and $p \in \Spec A$; set $B = A_p$, and
let $B^*$ be the completion of $B$. We prove that $B \to B^*$ is a regular
homomorphism; that is, for any prime ideal $p'$ of $B$, we need to show that
$B^* \otimes_B \kappa(p')$ is geometrically regular over $\kappa(p')$. However, $A/p' \cap A$ is also a
complete local ring, so that we can replace $A$ by $A/p' \cap A$ and reduce to the
case $p' = (0)$. Thus assume that $A$ is an integral domain, and let $L$ be the
common field of fractions of $A$ and $B$; we must show that $B^* \otimes_B L$ is
geometrically regular over $L$.

The problem can be further reduced to the case when $A$ is a regular local ring.
In fact, by Theorem~\hyperref[thm:29.5]{29.5}, $A$ contains a complete regular local ring $R$ and
is a finite module over $R$. Set $p \cap R = \mathfrak{q}$, $R_\mathfrak{q} = S$ and
$B' = A_p = A \otimes_R S$; then $B'$ is a semilocal ring, and $B$ is a localisation of $B'$
at a maximal ideal, so that $B^*$ is one direct factor of $B'^* = B' \otimes_S S^*$. Write $K$
for the common field of fractions of $R$ and $S$.
\[
\begin{tikzcd}
B'^* = B' \otimes_S S^* \arrow[r] & B^* \\
B' = A \otimes_R S \arrow[u] \arrow[r] & B = A_p \arrow[u] \arrow[r] & L \\
R \arrow[u] \arrow[r] & S = R_\mathfrak{q} \arrow[u] \arrow[r] & K \arrow[u]
\end{tikzcd}
\]
Now $B^* \otimes_B L$ can be written as $B^* \otimes_S L$, and it is hence a direct factor
of $B'^* \otimes_S L = S^* \otimes_S L = (S^* \otimes_S K) \otimes_K L$, so that we need only show
that $S^* \otimes_S K$ is geometrically regular over $K$.

Now $R$, $S$ and $S^*$ are regular local rings, and $S^* \otimes_S K$ is a localisation of
$S^*$, so is a regular ring. Hence if $\chark K = 0$, there is nothing to prove. We
assume that $\chark K = p$ in what follows. Then $R$ has a coefficient field $k$, and
can be written $R = k[[X_1, \ldots, X_n]]$. Choose a directed family $\{k_\alpha\}$ of cofinite
subfields $k_\alpha \subset k$ such that $\bigcap_\alpha k_\alpha = k^p$, and set
$R_\alpha = k_\alpha[[X_1^p, \ldots, X_n^p]]$; write $K_\alpha$ for the field of fractions of
$R_\alpha$. Then one sees easily (compare the proof of Theorem~\hyperref[thm:30.9]{30.9}) that
$\bigcap_\alpha K_\alpha = K^p$.

We set $\mathfrak{q}_\alpha = \mathfrak{q} \cap R_\alpha$; then since $R_\alpha \supset R^p$, we see that
$\mathfrak{q}$ is the unique prime ideal of $R$ lying over $\mathfrak{q}_\alpha$. Hence if we let
$S_\alpha = (R_\alpha)_{\mathfrak{q}_\alpha}$ then $S = R_\mathfrak{q} = R \otimes_{R_\alpha} S_\alpha$, and $S$ is a
finite module over $S_\alpha$. Hence $S^* = S \otimes_{S_\alpha} S_\alpha^*$; let us prove that $S^*$
is 0-smooth over $S$ relative to $S_\alpha$ (see \S\ref{sec:28}). Suppose we are given a
commutative diagram of the form
\[
\begin{tikzcd}
S^* \arrow[r, "\bar{u}"] & C/N \\
S_\alpha \arrow[u] \arrow[r] & S \arrow[u] \arrow[r] & C \arrow[ul]
\end{tikzcd}
\]
where $C$ is a ring and $N$ is an ideal of $C$ with $N^2 = 0$. If there is a lifting
$u'\colon S^* \to C$ of $\bar{u}$ as a homomorphism of $S_\alpha$-algebras, set
$w = u'\big|_{S^*}$, and let $\bar{v} = u \otimes w\colon S^* = S \otimes_{S_\alpha} S_\alpha^* \to C$; then
one sees easily that $\bar{v}$ is a lifting of $\bar{u}$ over $S$. Hence $S^*$ is 0-smooth over
$S$ relative to $S_\alpha$. Now for $Q \in \Spec(S^*)$, let $Q \cap S = (0)$; then $(S^*)_Q$ is
a local ring of $S^* \otimes_S K$, and conversely, every local ring of $S^* \otimes_S K$ is of
this form. From the diagram
\[
\begin{tikzcd}
S \arrow[r] & K \arrow[r] & (S^*)_Q/N(S^*)_Q \\
S_\alpha \arrow[u] \arrow[r] & K_\alpha \arrow[u] \arrow[ru, ""'] & (S^*)_Q \arrow[u]
\end{tikzcd}
\]
one sees that $(S^*)_Q$ is 0-smooth over $K$ relative to $K_\alpha$. Set $E = (S^*)_Q$ and
$m = \rad(E)$; then $E$ is $m$-smooth over $K$ relative to $K_\alpha$, so that according
to Theorem~\hyperref[thm:28.4]{28.4},
\[
  \Omega_{K/K_\alpha} \otimes_K (E/m) \to \Omega_{E/K} \otimes_E (E/m)
\]
is injective for every $\alpha$. Moreover, since $\bigcap_\alpha K_\alpha = K^p$, by
\S\ref{sec:30}, Lemma~\hyperref[lem:30.4]{4},
\[
  \Omega_K \to \varprojlim \Omega_{K/K_\alpha}
\]
is injective, and hence
\[
  \Omega_K \otimes (E/m) \to \varprojlim (\Omega_{K/K_\alpha} \otimes (E/m))
\]
is also injective. Therefore, from the commutative diagram
\[
\begin{tikzcd}
\Omega_K \otimes (E/m) \arrow[r] & \Omega_E \otimes (E/m) \\
\varprojlim (\Omega_{K/K_\alpha} \otimes (E/m)) \arrow[u] \arrow[r] &
\varprojlim (\Omega_{E/K_\alpha} \otimes (E/m)) \arrow[u]
\end{tikzcd}
\]
we finally see that $\Omega_K \otimes (E/m) \to \Omega_E \otimes (E/m)$ is injective. Hence it
follows from the corollary to Theorem~\hyperref[thm:28.6]{28.6} that $E$ is $m$-smooth over $K$, and
thus is geometrically regular. Since $E$ is an arbitrary local ring of $S^* \otimes_S K$,
we see that $S^* \otimes_S K$ is geometrically regular over $K$.
\end{proof}

\begin{theorem}\label{thm:32.4}
Let $A$ be a Noetherian ring; if $A_m \to (A_m)^*$ is regular for every maximal ideal
$m$ of $A$, then $A$ is a G-ring.
\end{theorem}

\begin{proof}
Since $(A_m)^*$ is a G-ring, by Theorem~\hyperref[thm:32.2]{2}, $A_m$ is also a G-ring. For any
$p \in \Spec A$, if we let $m$ be a maximal ideal of $A$ containing $p$ then $A_p$ is a
localisation of the G-ring $A_m$, and hence $A_p \to (A_p)^*$ is regular.
\end{proof}

Theorem~\hyperref[thm:32.4]{4} makes it much easier to distinguish G-rings. For example, the
next theorem is based on Theorem~\hyperref[thm:32.4]{4}.

\begin{theorem}\label{thm:32.5}
Let $A$ be a Noetherian semilocal ring; then a sufficient condition for $A$ to be
a G-ring is that if $C$ is a finite $A$-algebra which is an integral domain, $m$ is
a maximal ideal of $C$ and we write $B = C_m$, then $(B^*)_Q$ is a regular local
ring for every $Q \in \Spec(B^*)$ such that $Q \cap B = (0)$.
\end{theorem}

\begin{remark}
It is easy to see that the condition is also necessary.
\end{remark}

\begin{proof}
By the previous theorem, we need only show that under the given condition,
$A \to A^*$ is regular. Let $p \in \Spec A$, and let $L$ be a finite extension of
$\kappa(p)$; we prove that $A^* \otimes_A L$ is regular. Suppose that
$L = \kappa(p)(\xi_1, \ldots, \xi_m)$; then multiplying each $\xi_i$ by an element of $A/p$ we
can assume that $\xi_i$ is integral over $A/p$, so that if we set
$C = (A/p)[\xi_1, \ldots, \xi_m]$, then $C$ is a finite $A$-module, and the field of
fractions of $C$ is $L$. Now $C^* = A^* \otimes_A C$, and if we write
$m_1, \ldots, m_s$ for the maximal ideals of $C$ and set $B_i = C_{m_i}$, then
$C^* = B_1^* \times \cdots \times B_s^*$. We can identify any local ring of
$A^* \otimes_A L = C^* \otimes_C L$ with the localisation $(B_i^*)_Q$ of one of the factors
$B_i^*$ at some prime ideal $Q$ of $B_i^*$ with $Q \cap B_i = (0)$, and by assumption
this is regular. Hence $A^* \otimes_A L$ is a regular ring.
\end{proof}

\begin{theorem}[\textbf{H.~Mizutani}]\label{thm:32.6}
Let $R$ be a regular ring. If the weak Jacobian condition (WJ) of \S\ref{sec:30} holds for
$R[X_1, \ldots, X_n]$ for every $n \ge 0$, then $R$ is a G-ring.
\end{theorem}

\begin{proof}
Since (WJ) is inherited by any localisation we can assume that $R$ is local. We
prove that the condition of Theorem~\hyperref[thm:32.5]{5} holds. Set
$R_n = R[X_1, \ldots, X_n]$. Any integral domain $C$ which is finite as an $R$-module
can be expressed as $C = R_n/Q$ with $Q \in \Spec(R_n)$ for some $n$. Let $m$ be a
maximal ideal of $C$, $M$ the maximal ideal of $R_n$ corresponding to $m$, and
$S = (R_n)_M$; then it is enough to show that $(S^*)_P/Q(S^*)_P$ is regular for every
$P \in \Spec(S^*)$ with $P \cap S = QS$. If $\htop Q = r$ then $\htop Q(S^*)_P = r$, and by
assumption there exist $D_1, \ldots, D_r \in \Der(R_n)$ and
$f_1, \ldots, f_r \in Q$ such that $\det(D_if_j) \notin Q$. Now $D_i$ has a natural extension
to $S$, and then to $S^*$, and since $P \cap R_n = Q$, we have
$\det(D_if_j) \notin P$, so that by Theorem~\hyperref[thm:30.4]{30.4}, $(S^*)_P/Q(S^*)_P$ is
regular.
\end{proof}

\begin{corollary*}
A ring which is finitely generated over a field, or a localisation of such a
ring, is a G-ring.
\end{corollary*}

\begin{proof}
It follows from the definition that a quotient or localisation of a G-ring is
again a G-ring, so that we need only show that for a field $k$, the ring
$k[X_1, \ldots, X_n]$ is a G-ring; but by Theorems~\hyperref[thm:30.3]{30.3} and \hyperref[thm:30.5]{30.5}, (WJ)
holds in $k[X_1, \ldots, X_{n+m}]$, so that $k[X_1, \ldots, X_n]$ is a G-ring by the
theorem.
\end{proof}

\begin{remark}
The local rings which appear in algebraic geometry are essentially of finite
type over a field, and therefore G-rings. Hence for these rings, properties such
as reduced and normal pass to the completion.
\end{remark}

\begin{remark}
If $R$ is a regular ring containing a field of characteristic $0$, and (WJ) holds
in $R$, then by Ex.~\hyperref[ex:30.4-ex]{30.4}, it also holds in $R[X_1, \ldots, X_n]$. Hence $R$ and
$R[X_1, \ldots, X_n]$ are G-rings. In particular by Theorem~\hyperref[thm:30.8]{30.8}, rings of
convergent power series over $\mathbb{R}$ and $\mathbb{C}$ are G-rings.
\end{remark}

A theorem proved by Grothendieck asserts that if $A$ is a G-ring, then so is
$A[X]$; the proof is very hard, and we omit it, referring only to [M], Theorem 77.
The analogous statement for $A[[X]]$ remained unsolved for a long time, but was
recently proved for a semilocal ring $A$ by C.~Rotthaus [3]. In the non-semilocal
case she proved in [4] that, if $A$ is a finite-dimensional excellent ring
containing the rational numbers, then $A[[X]]$ is excellent. On the other hand,
Nishimura [3] showed that there exists a G-ring $A$ such that $A[[X]]$ is not a
G-ring.

Nagata [8] studied the condition that $\Reg(A)$ is open in $\Spec A$; putting
together Nagata's work with his own theory of G-rings, Grothendieck gave the
definition of excellent ring in [G2].

\begin{definition*}
A Noetherian ring $A$ is \emph{excellent} if it satisfies the following three conditions:
\begin{enumerate}[(1)]
  \item $A$ is universally catenary;
  \item $A$ is a G-ring;
  \item $\Reg(B) \subset \Spec B$ is open for every finitely generated $A$-algebra $B$.
\end{enumerate}
A Noetherian ring satisfying (2) and (3) is said to be \emph{quasi-excellent}.
\end{definition*}

One can prove that the classes of rings satisfying each of (1), (2) and (3) are
closed under localisation, finitely generated extensions and passing to
quotients. It can also be proved that (2) implies (3) for semilocal Noetherian
rings. A complete Noetherian local ring is excellent, as are practically all
Noetherian rings in applications. For more information on excellent rings, see
[M], Ch.~13, or [G2].

R.~Y.~Sharp [S] defined the notion of an \emph{acceptable ring}, replacing condition
(2) by the condition that all formal fibres of all localisations of $A$ are
Gorenstein, and replacing $\Reg(B)$ by $\Gor(B)$ in (3), and showed that the
resulting theory is analogous to the theory of excellent rings (see also
Greco-Marinari [1], Sharp [6]).

Using his cohomology theory, M.~Andr\'e [1] proved the following theorem. Let
$A$, $B$ be Noetherian local rings, and $\varphi\colon A \to B$ a local homomorphism;
suppose that $A$ is quasi-excellent and $B$ is $m_B$-smooth over $A$, where
$m_B = \rad(B)$; then $\varphi$ is regular. This is an extremely strong theorem; the
result of Rotthaus mentioned above also makes use of this.
